\section{Original section rings inside the original scheme function field}

\begin{definition}[The actual generic-germ algebra map]\label{def:sc0-scheme-sections-function-field-map}
\lean{AlgebraicGeometry.Scheme.Hom.sectionsFunctionFieldAlgHom}
\leanok
For an integral scheme $X$ with its original structure morphism
$f:X\to\operatorname{Spec} k$ and a nonempty original open $U$, use
the actual generic germ $\Gamma(X,U)\to K(X)$. This is a $k$-algebra
map for the original section and generic-stalk algebra structures.
\end{definition}
\begin{proof}
The generic point belongs to every nonempty open of the irreducible
space. The actual germ is a ring map. Its composite with the original
base-field section map is exactly the original base-field stalk map,
by the proved structure-morphism germ compatibility, proving the
algebra-map assertion.
\end{proof}

\begin{definition}[Literal section subalgebra]\label{def:sc0-scheme-sections-function-field-subalgebra}
\lean{AlgebraicGeometry.Scheme.Hom.sectionsFunctionFieldSubalgebra}
\leanok
\uses{def:sc0-scheme-sections-function-field-map}
Let $C_U\subset K(X)$ be the actual range subalgebra of this generic germ.
\end{definition}

\begin{theorem}[The actual original section-ring equivalence]\label{thm:sc0-scheme-sections-function-field-equiv}
\lean{AlgebraicGeometry.Scheme.Hom.sectionsFunctionFieldEquiv}
\leanok
\uses{def:sc0-scheme-sections-function-field-subalgebra}
The actual generic germ induces a $k$-algebra equivalence
$\Gamma(X,U)\simeq C_U$.
\end{theorem}
\begin{proof}
On an integral scheme the germ map from any nonempty open to any
of its stalks is injective. In particular the actual generic germ is
injective, and its range is exactly $C_U$ by definition. The actual
injective-algebra-map range equivalence therefore gives the assertion.
\end{proof}

\begin{theorem}[The actual value of the equivalence]\label{thm:sc0-scheme-sections-function-field-equiv-apply}
\lean{AlgebraicGeometry.Scheme.Hom.sectionsFunctionFieldEquiv_apply}
\leanok
\uses{thm:sc0-scheme-sections-function-field-equiv}
For every original section $s$, its image under this equivalence,
viewed in $K(X)$, is its actual generic germ.
\end{theorem}
\begin{proof}
The actual range equivalence is defined by the original algebra map.
\end{proof}

\begin{theorem}[Restriction gives literal field containment]\label{thm:sc0-scheme-sections-function-field-containment}
\lean{AlgebraicGeometry.Scheme.Hom.sectionsFunctionFieldSubalgebra_le}
\leanok
\uses{def:sc0-scheme-sections-function-field-subalgebra}
If $U\subseteq V$ are nonempty original opens, then $C_V\subseteq C_U$
as literal subalgebras of the same actual function field $K(X)$.
\end{theorem}
\begin{proof}
Represent an arbitrary element of $C_V$ by an actual section over
$V$. Restrict that section to $U$. Compatibility of restriction with
the actual generic germ shows that its value in $K(X)$ is unchanged,
so the element belongs to the actual range over $U$.
\end{proof}

\begin{theorem}[Actual restriction square]\label{thm:sc0-scheme-sections-function-field-restriction}
\lean{AlgebraicGeometry.Scheme.Hom.sectionsFunctionFieldEquiv_restriction}
\leanok
\uses{thm:sc0-scheme-sections-function-field-equiv,thm:sc0-scheme-sections-function-field-containment}
The original section restriction from $V$ to $U$ corresponds under
the actual equivalences to the literal inclusion $C_V\subseteq C_U$.
\end{theorem}
\begin{proof}
View both sides in the actual function field. The preceding explicit
value formula reduces their equality to actual generic-germ
compatibility with the original scheme restriction.
\end{proof}

\begin{theorem}[The original function field remains the actual fraction field]\label{thm:sc0-scheme-sections-function-field-fraction}
\lean{AlgebraicGeometry.Scheme.Hom.sectionsFunctionFieldSubalgebra_isFractionRing}
\leanok
\uses{thm:sc0-scheme-sections-function-field-equiv}
For an original nonempty affine open $U$, the actual $K(X)$ is the
fraction field of the literal subalgebra $C_U\subset K(X)$.
\end{theorem}
\begin{proof}
The actual generic stalk of an integral scheme is the fraction field
of each nonempty affine section ring. Transport this proved
localization property along the actual section-ring equivalence.
The explicit value of that equivalence shows that its composite into
the actual function field is precisely the original generic germ,
so this is the canonical literal subalgebra algebra structure.
\end{proof}

\begin{theorem}[Actual principal localization for intrinsic basic opens]\label{thm:sc0-scheme-sections-function-field-basic-localization}
\lean{AlgebraicGeometry.Scheme.Hom.sectionsFunctionFieldSubalgebra_isLocalization_basicOpen}
\leanok
\uses{thm:sc0-scheme-sections-function-field-restriction}
For an actual nonempty basic open $D(s)$ of a nonempty affine $U$,
the literal inclusion $C_U\subset C_{D(s)}$ is principal localization
at the actual image of the original section $s$.
\end{theorem}
\begin{proof}
The actual original restriction $\Gamma(X,U)\to\Gamma(X,D(s))$ is
localization away from $s$. Use the actual equivalences on both
section rings and the proved restriction square. The image of $s$
is a unit because its original restriction is a unit. Every element
of $C_{D(s)}$ is the actual image of an original localized section,
which has an actual numerator and a denominator power of $s$;
applying the target equivalence to this fraction equation gives
the required numerator and denominator in $C_U$. Finally the literal
inclusion is injective, so equality descends with denominator one.
These are the three actual localization axioms.
\end{proof}
